\chapter{پژوهش‌های مرتبط}
\label{ch:literature}

\section{چارچوب مقایسه روش‌ها}
روش‌های آموزش هم‌آورد سریع را می‌توان بر پایه نقطه‌ای که در زنجیره آموزش تغییر می‌دهند مقایسه کرد. برخی روش‌ها توزیع آغاز حمله را عوض می‌کنند، برخی شکل تابع هزینه را با منظم‌سازی محدود می‌کنند و گروهی اندازه گام یا اطلاعات انتقال‌یافته از گذشته را تغییر می‌دهند. این دسته‌بندی کاملاً جدا از هم نیست؛ برای نمونه، یک روش می‌تواند هم آغاز تصادفی داشته باشد و هم از منظم‌سازی استفاده کند.

در این فصل، معیار اصلی مقایسه تنها بیشترین درستی روی یک دادگان نیست. یک روش سریع باید کیفیت حل درونی، هزینه گذرهای اضافی، حافظه، حساسیت ابرپارامتر و قابلیت استفاده در تنظیم‌های مختلف را نیز در نظر بگیرد. این نگاه با هدف پیشنهاد پایان‌نامه، یعنی پیش‌بینی کم‌هزینه بیش‌برازش، سازگار است. بخش مهمی از مرور زیر در سمینار پیشین نیز بررسی شده بود؛ در اینجا ارتباط آن با نتایج نهایی \sora{} روشن‌تر می‌شود.

\section{آموزش چندگامی و بده‌بستان ستبری}
آموزش مبتنی بر \lr{PGD} مسئله درونی را با چند حرکت در جهت افزایش هزینه حل می‌کند \cite{madry2018towards}. مزیت آن این است که پس از هر حرکت، شیب دوباره محاسبه می‌شود؛ بنابراین، الگوریتم می‌تواند با تغییر جهت محلی تابع هزینه سازگار شود. این انعطاف در برابر هندسه غیرخطی، یکی از دلایلی است که ارزیابی چندگامی می‌تواند شکست پنهان‌شده در حمله تک‌گامی را آشکار کند. هزینه آن نیز از همین محاسبات تکراری ناشی می‌شود.

روش \lr{TRADES} رابطه میان خطای طبیعی و خطای مرزی را با یک هدف آموزشی متشکل از هزینه پاک و واگرایی خروجی پاک و آشفته بررسی می‌کند \cite{zhang2019theoretically}. یک صورت متداول هدف آن چنین است:
\begin{equation}
\mathbb E_{(x,y)}\left[\mathcal L(f_\theta(x),y)+\lambda
\max_{\delta\in\Delta_\infty(x,\epsilon)}
D_{\mathrm{KL}}\bigl(q_\theta(\cdot\mid x)\|q_\theta(\cdot\mid x+\delta)\bigr)\right].
\end{equation}
ضریب $\lambda$ در این نوشتار وزن جمله ستبری است؛ در برخی منابع از معکوس آن استفاده می‌شود. مقایسه تنظیم‌ها بدون توجه به این قرارداد می‌تواند گمراه‌کننده باشد. این روش در مقاله \sora{} در کنار \lr{PGD-2} و \lr{PGD-10} به‌عنوان مرجع چندگامی به کار رفته است.

پژوهش‌های مربوط به بده‌بستان درستی پاک و ستبری نشان می‌دهند افزایش یکی الزاماً بدون هزینه برای دیگری نیست \cite{tsipras2019robustness}. بنابراین، برتری یک دفاع باید با هر دو معیار سنجیده شود. از سوی دیگر، این بده‌بستان یک دلیل برای پذیرفتن هر مقدار افت پاک نیست؛ مقایسه روش‌ها می‌تواند نشان دهد یک سازوکار کاراتر، برای سطح مشابه ستبری درستی پاک بیشتری حفظ می‌کند.

\section{استفاده مجدد از محاسبات}
\subsection{آموزش هم‌آورد رایگان}
روش \lr{Free AT} از یک گذر پس‌انتشار برای به‌روزرسانی هم‌زمان آشفتگی و پارامترهای شبکه استفاده می‌کند \cite{shafahi2019adversarial}. برای مؤثربودن این استفاده مشترک، دسته آموزشی چند بار بازپخش می‌شود. این راهبرد هزینه تولید حمله را با آموزش وزن‌ها درهم می‌آمیزد و نشان می‌دهد تفکیک کامل «محاسبه حمله» از «محاسبه آموزش» همواره ضروری نیست.

واژه رایگان در این زمینه نسبی است. بازپخش دسته همچنان محاسبه دارد و مقایسه تعداد دوره‌های اسمی بدون دانستن تعداد بازپخش‌ها کافی نیست. در مطالعه \sora{}، این روش با تنظیمی متمایز از بسیاری از روش‌های دیگر گزارش شده است. ازاین‌رو، هم تعداد گذر و هم زمان واقعی باید با درنظرگرفتن پروتکل کامل تفسیر شوند.

\subsection{یک گام همراه با آغاز تصادفی}
وونگ و همکاران نشان دادند آغاز تصادفی می‌تواند آموزش تک‌گامی را به‌طور چشمگیری بهبود دهد \cite{wong2020fast}. نمونه ابتدا در همسایگی تصویر جابه‌جا می‌شود و سپس یک گام علامت شیب اعمال می‌شود. این تصادفی‌سازی، وابستگی مولد حمله به یک نقطه آغاز ثابت را کاهش می‌دهد. بااین‌حال، همان پژوهش مسئله بیش‌برازش فاجعه‌بار را نیز برجسته کرد.

نکته مرتبط با پایان‌نامه آن است که تصادفی‌کردن نقطه آغاز الزاماً تمام اندازه‌ها و جهت‌های مؤثر را با کیفیت یکسان پوشش نمی‌دهد. تصویرشدن به کره و یک گام نسبتاً بزرگ می‌تواند بخش زیادی از نمونه‌ها را به مرز ببرد. در نتیجه، توزیع نهایی آشفتگی حاصل ترکیب نویز، جهت شیب و برش است؛ تنها دانستن توزیع نویز برای توصیف آن کافی نیست.

\section{روش‌های مبتنی بر تنوع آشفتگی}
روش \lr{N-FGSM} دامنه نویز آغازین را افزایش می‌دهد و از برش نهایی آشفتگی به کره آموزشی صرف‌نظر می‌کند \cite{dejorge2022make}. هدف آن افزایش تنوع و کاهش شکست آموزش تک‌گامی است. این روش از نظر تعداد گذرهای اصلی در گروه روش‌های واقعاً سریع قرار می‌گیرد و در بسیاری از معیارهای طبیعی عملکرد رقابتی دارد.

این نتیجه به معنای یکسان‌بودن دامنه آشفتگی آموزش و ارزیابی نیست. استفاده از ورودی‌های آموزشی بیرون از کره ارزیابی یک انتخاب طراحی است و باید هنگام توضیح روش تصریح شود. همچنین، افزایش نویز ممکن است در یک دادگان به پایداری کمک کند و در دادگان دیگر درستی پاک را کاهش دهد. این مسئله انگیزه‌ای برای انتخاب وابسته به وضعیت مدل، به‌جای اتکا به یک مقدار ثابت بزرگ، فراهم می‌کند.

از دیدگاه بیش‌برازش اپسیلون، تصادفی‌سازی می‌تواند احتمال تمرکز مدل بر چند فاصله محدود را کاهش دهد. بااین‌حال، این تفسیر یک توضیح سازوکاری است و نه اثبات اینکه هر نوع نویز مستقل از توزیع داده، بیش‌برازش را حذف می‌کند. نتایج گسترده مقاله نشان می‌دهد ارزیابی خارج از \lr{CIFAR-10} برای تشخیص این تفاوت‌ها ضروری است \cite{sora2026}.

\section{منظم‌سازی هندسه محلی}
\subsection{هم‌راستایی شیب}
روش \lr{GradAlign} از کاهش هم‌راستایی شیب در همسایگی ورودی به‌عنوان نشانه نامطلوب استفاده می‌کند و آن را در تابع هدف جریمه می‌کند \cite{andriushchenko2020understanding}. شکل مفهومی جمله منظم‌ساز عبارت است از
\begin{equation}
R_{\mathrm{GA}}=\mathbb E_{\eta}\left[1-
\cos\bigl(\nabla_x\ell(x),\nabla_x\ell(x+\eta)\bigr)\right].
\end{equation}
کمینه‌کردن این جمله، مدل را به حفظ جهت شیب در یک همسایگی تصادفی تشویق می‌کند. این انتخاب به‌طور مستقیم با ایده محلی‌بودن تقریب تک‌گامی ارتباط دارد.

تفاوت آن با \pertalign{} در نقش و نقاط مقایسه است. \lr{GradAlign} یک قید آموزشی است و برای مشتق‌گیری جمله منظم‌ساز نسبت به وزن‌ها، محاسبات بیشتری لازم می‌شود. \pertalign{} شیب ابتدا و انتهای همان حمله جاری را مقایسه می‌کند و نخست به‌عنوان ابزار پایش تعریف می‌شود. دو شاخص ممکن است رفتار هشدار مشابهی داشته باشند، اما یکسان‌بودن مقدار یا هزینه آن‌ها از این شباهت نتیجه نمی‌شود.

\subsection{هموارسازی خروجی و خطی‌بودن محلی}
روش \lr{NuAT} از منظم‌سازی مبتنی بر نرم هسته‌ای برای بهبود رفتار محلی خروجی بهره می‌گیرد \cite{sriramanan2021towards}. هدف کلی این خانواده آن است که آموزش تنها به نقطه انتهایی حمله محدود نماند و نوعی سازگاری در اطراف داده ایجاد شود. بااین‌حال، جمله اضافی تابع هدف می‌تواند به گذر پیشرو، ذخیره‌سازی یا مشتق‌گیری بیشتری نیاز داشته باشد.

روش \lr{ELLE} بر سنجش و منظم‌سازی خطای تقریب خطی تابع هزینه تمرکز می‌کند \cite{rocamora2024efficient}. این دیدگاه به مسئله‌ای نزدیک است که در \lr{FGSM} رخ می‌دهد: اگر تقریب خطی محلی معتبر نباشد، اندازه و جهت حرکت حاصل از آن نیز ممکن است نامناسب شود. مزیت روش، ارتباط روشن با هندسه است؛ محدودیت عملی آن باید با هزینه واقعی و حساسیت تنظیم‌ها سنجیده شود، نه فقط با تعداد گام اسمی حمله.

\subsection{نمونه‌های هم‌آورد غیرعادی}
در روش \lr{AAER}، نمونه‌ای که پس از حمله هزینه کمتری از ورودی آغازین دارد، نمونه هم‌آورد غیرعادی در نظر گرفته می‌شود \cite{lin2023eliminating}. وجود چنین نمونه‌ای نشان می‌دهد حمله با هدف افزایش هزینه، دست‌کم در نقطه انتهایی به نتیجه مورد انتظار نرسیده است. فراوانی آن‌ها می‌تواند با بیش‌برازش فاجعه‌بار افزایش یابد.

این شاخص از نظر مفهومی مستقیم است، اما زمان ظهور آن اهمیت دارد. ممکن است جهت شیب پیش از آنکه هزینه انتهایی از هزینه اولیه کمتر شود، به‌شدت تغییر کند. در این حالت، یک سنجه شیب‌محور می‌تواند زودتر واکنش نشان دهد. مقاله \sora{} این تفاوت را در نمودارهای زمانی مقایسه می‌کند، ولی برای نتیجه‌گیری درباره نرخ عمومی هشدار، ارزیابی آماری جداگانه لازم است.

\section{انتخاب تطبیقی اندازه گام}
ثابت‌گرفتن اندازه گام فرض می‌کند هندسه مناسب حمله در طول آموزش تغییر اساسی نمی‌کند. روش‌های تطبیقی این فرض را تضعیف می‌کنند و از نشانه‌های مدل برای کنترل حرکت بهره می‌برند. \lr{ATAS} اندازه گام را با اطلاعات بزرگی شیب تنظیم می‌کند \cite{huang2023fast}. این راهبرد می‌تواند از حرکت نامناسب در نمونه‌هایی با ویژگی‌های متفاوت بکاهد، اما بزرگی شیب و خمیدگی جهتی دو کمیت متفاوت‌اند.

روش‌های دیگری از یادگیری تقویتی برای تعیین آشفتگی یا از شباهت نویز و شیب برای تنظیم گام استفاده کرده‌اند \cite{nie2021adaptive,zhao2025avoiding}. مزیت عمومی آن‌ها، واکنش به وضعیت جاری است. در مقابل، هر سازوکار تطبیقی ممکن است خودش به انتخاب آستانه، دامنه عمل یا پارامترهای کنترلی نیاز داشته باشد. بنابراین، «تطبیقی» بودن به‌تنهایی به معنای بی‌نیازی از تنظیم نیست.

تمایز \sora{} آن است که قانون کنترل گام از تقریب درجه دوم هزینه در جهت حرکت استخراج می‌شود. نسبت میان شیب آغاز و شیب پایان، تخمینی از خمیدگی جهتی فراهم می‌کند. سپس برای استفاده عملی، این تخمین هموار و با تأخیر به دسته بعد منتقل می‌شود. نسخه عملی، به دلیل ثابت‌کردن صورت کسر و تصادفی‌کردن گام مؤلفه‌ها، عین حل دقیق یک مسئله درجه دوم ثابت نیست؛ این تمایز در فصل روش حفظ می‌شود.

\section{حذف مؤلفه‌های نامطمئن و استفاده از سابقه}
روش‌های \lr{ZeroGrad} و \lr{MultiGrad} با حذف یا بررسی پایداری مؤلفه‌های شیب می‌کوشند از اثر جهت‌های نامطمئن بکاهند \cite{golgooni2023zerograd}. در یک تابع تقریباً خطی، انتظار می‌رود علامت شیب در همسایگی پایدارتر باشد. مؤلفه‌های کوچک‌تر نیز ممکن است نسبت به تغییر محلی حساس‌تر باشند؛ بنابراین، کنارگذاشتن بخشی از آن‌ها می‌تواند نقش یک اصلاح کم‌هزینه را ایفا کند.

بااین‌حال، حذف مؤلفه با تنظیم بزرگی حرکت در تمام مؤلفه‌ها یکسان نیست. همچنین، بررسی چند آغاز یا چند شیب ممکن است هزینه را افزایش دهد. مقاله \sora{} این روش‌ها را هم از نظر ستبری و هم از نظر زمان و حافظه مقایسه می‌کند. نتیجه یک تنظیم خاص نباید به همه ترکیب‌های دادگان و معماری تعمیم داده شود.

\section{اطلاعات مرتبه دوم در دفاع هم‌آورد}
رابطه میان خمیدگی مرز تصمیم و آسیب‌پذیری، انگیزه روش‌های منظم‌سازی خمیدگی بوده است \cite{moosavi2019robustness}. این خانواده به‌جای کنترل تنها یک حمله، تلاش می‌کند هندسه مدل را تغییر دهد. رویکردهای دیگر نیز از تقریب‌های هسین یا اطلاعات مرتبه دوم برای تولید حمله یا دفاع استفاده می‌کنند. چالش اصلی برای آموزش سریع، به‌دست‌آوردن این اطلاعات با هزینه مناسب است.

دو نکته باید تفکیک شوند. نخست، کم‌بودن تغییر جهت شیب در یک مسیر، به معنای کوچک‌بودن تمام مقادیر ویژه هسین نیست. دوم، کوچک‌بودن یک مؤلفه خمیدگی، به‌تنهایی ستبری دسته‌بند را تضمین نمی‌کند؛ فاصله از مرز تصمیم و بزرگی شیب نیز مؤثرند. تحلیل \pertalign{} در همین دامنه محلی و جهت‌وابسته تفسیر می‌شود.

\section{جایگاه پژوهش حاضر}
پژوهش حاضر از سه ایده شناخته‌شده بهره می‌گیرد: اهمیت هندسه محلی، سودمندی تنوع آشفتگی و امکان استفاده مجدد از محاسبات آموزش. پیوند این ایده‌ها در \sora{} چنین است: دو شیب لازم برای تولید نمونه و آموزش وزن‌ها، برای پایش تغییر جهت و تخمین خمیدگی جهتی نیز استفاده می‌شوند؛ اندازه گام به وضعیت مدل واکنش نشان می‌دهد؛ و نمونه‌برداری مؤلفه‌ای از تمرکز بر یک اندازه ثابت می‌کاهد.

% Source: Newly compiled qualitative comparison from seminar literature review
% and SORA Section 2; this exact table was not in the paper or seminar.
\begin{table}[!htbp]
\centering
\caption[مقایسه سازوکارهای آموزش هم‌آورد سریع]{مقایسه سازوکارهای مرتبط؛ تلفیق مرور سمینار و مقاله، بدون ادعای برتری عددی.}
\label{tab:related-mechanisms}
\small
\begin{tabular}{p{.19\textwidth}p{.34\textwidth}p{.34\textwidth}}
\toprule
روش & اطلاعات یا تغییر اصلی & نکته مرتبط با این پژوهش\\
\midrule
\lr{FGSM-RS} & آغاز تصادفی و یک گام & وابستگی کیفیت به توزیع نهایی\\
\lr{N-FGSM} & نویز بزرگ‌تر، حذف برش به کره & دامنه آموزشی متفاوت از ارزیابی\\
\lr{GradAlign} & منظم‌سازی هم‌راستایی شیب & کنترل هندسه با محاسبه اضافی\\
\lr{AAER} & کنترل نمونه‌های غیرعادی & بررسی نقطه انتهایی حمله\\
\lr{ELLE} & منظم‌سازی خطی‌بودن محلی & کنترل خطای تقریب مرتبه اول\\
\lr{ATAS} & گام وابسته به بزرگی شیب & تطبیق با اطلاعات مرتبه اول\\
\sora & دو شیب موجود و گام تطبیقی & تخمین خمیدگی جهتی با تأخیر\\
\bottomrule
\end{tabular}
\end{table}

جدول \ref{tab:related-mechanisms} خلاصه‌ای مفهومی ارائه می‌کند. قضاوت نهایی درباره کارایی باید بر نتایج فصل آزمایش‌ها متکی باشد. به‌ویژه، ثابت‌بودن ابرپارامترهای اختصاصی \sora{} در مجموعه ارزیابی‌شده، ادعایی تجربی درباره همان تنظیم‌هاست و تضمین نمی‌کند هر دادگان تازه نیز بدون تنظیم به همان کیفیت برسد.
